\subsection{Ordinary eighth moments and the nonunit-pair outlet}\label{proof:R2-C046}
\paragraph{R46.1 Local statements and conventions}
Use the original finite probability space and operators of
Section~\ref{sec:prelim}. A terminal residual has at least two distinct
lengths. An ordinary component consists of that terminal, its connector
of length \(L\ge1\), one existing receiving arm of length \(m\ge3\),
and the ordered negative post block specified in
\mbox{\hyperref[part:R46.3]{part R46.3}}. The message includes its
root and terminal equal-pair factors. The directions in that directory
satisfy, for every finite host,
\fitdisplay{
E H^8\le e^{-d/1024},\qquad d=\|K1\|^2.
}{R46.1.2}

For a root whose strict residual is \((s,m)\), \(2\le s<m\), every
nonreducible terminal and every pair of independent positive lengths
\(L,M\) admit the output constructed below, with
\fitdisplay{
I_q(U)^3\le e^{-\eta_{\rm pair}d},
 \qquad \eta_{\rm pair}=\frac1{72}+\frac1{1024}
                      =\frac{137}{9216}.
}{R46.1.3}

\paragraph{R46.2 Equal pairs}
Equal-pair factors satisfy \eqref{q3:eq:equalpair} and its complete
reversal. They may be restored inside their positive terminal connector
or at the receiving root, where they only decrease the message.
Unequal post pairs retain their displayed directions and ordering.


\paragraph{R46.3 Exhaustive actual ordinary-eighth directory}\label{part:R46.3}

Let \(m\ge 3\) be one existing receiving arm, \(L\ge 1\) its original connector
to a terminal with strict residual size \(\ge2\), and

\fitdisplay{
P=\prod_{\nu=1}^k W_{u_\nu}\bar f_{v_\nu},\qquad
 4\le u_1<v_1<\cdots<u_k<v_k.
}{R46.3.1}

Impose \(m<u_{1}\) when \(k\ge 1.\) The restriction is imposed uniformly, even in rows where a weaker one would suffice. There
is no restriction between the terminal lengths and the other blocks.
Every post is an actual original pair of unequal arms.

The complete message is \(H=E_rP\) times the displayed anchor/connector
function, with the terminal equal pairs restored inside its actual
positive connector. Empty products equal one. Canonical pairs mean
\(h_a f_b\) for adjacent increasing residual lengths; their negatives
are \(W_a \bar f_b.\)

\paragraph{R46.3.1 Even terminal lists}

Write the strict terminal list as \(a_1<b_1<\cdots<a_j<b_j.\)


\begingroup\small
\setlength{\tabcolsep}{3pt}
\setlength{\LTleft}{0pt}\setlength{\LTright}{\fill}
\begin{longtable}{@{}p{0.2200\linewidth}p{0.2333\linewidth}p{0.2333\linewidth}p{0.2333\linewidth}@{}}
\caption{Ordinary eighth moments and the nonunit-pair outlet: parameter values.}\label{tab:parameters:r2-c046:1}\\
\toprule
Entry & L/anchor condition & Unpaired anchor/connector factor \(H_0\) & Eighth loss \\
\midrule\endfirsthead
\toprule
Entry & L/anchor condition & Unpaired anchor/connector factor \(H_0\) & Eighth loss \\
\midrule\endhead
\bottomrule\endfoot
\(j=1,\allowbreak a_{1}=1\) & \(L=1\); & \(f_m T(p \bar f_{b_1})\); & \(1/1024\). \\
\addlinespace[2pt]
\(j=1,\allowbreak a_{1}=1\) & \(L\ge 2\); & \(W_m T^L(q f_{b_1})\); & \(3/20\). \\
\addlinespace[2pt]
\(j=1,\allowbreak a_{1}=2\) & \(L=1\); & \(f_m T(f_{2} \bar f_{b_1})\); & at least \(1/4\). \\
\addlinespace[2pt]
\(j=1,\allowbreak a_{1}=2\) & \(L\ge 2\); & \(W_m T^L(\bar f_2 W_{b_1})\); & at least \(1/4\). \\
\addlinespace[2pt]
\(j=1,\allowbreak a_{1}=3\) & \(L=1\); & \(f_m T(f_{3} h_{b_1})\); & at least \(1/4\). \\
\addlinespace[2pt]
\(j=1,\allowbreak a_{1}=3\) & \(L\ge 2\); & \(W_m T^L(f_{3} h_{b_1})\); & at least \(1/4\). \\
\addlinespace[2pt]
\(j\ge 2,\allowbreak a_{1}=1\) & \(L=1\); & \(f_m T(\overline{F}),\allowbreak  F=q f_{b_1} \prod_{i\ge 2}h_{a_i} f_{b_i}\); & at least \(8/7+1/1024\). \\
\addlinespace[2pt]
\(j\ge 2,\allowbreak a_{1}=1\) & \(L\ge 2\); & \(W_m T^L(F),\allowbreak \) same \(F\); & at least \(8/7+1/1024\). \\
\addlinespace[2pt]
\(j\ge 2,\allowbreak a_{1}=2\) & \(L=1\); & \(\bar f_m A(F),\allowbreak  F=\prod_i h_{a_i} f_{b_i}\); & 2. \\
\addlinespace[2pt]
\(j\ge 2,\allowbreak a_{1}=2\) & \(L\ge 2\); & \(W_m T^L(F),\allowbreak \) same \(F\); & 2. \\
\addlinespace[2pt]
\(j\ge 2,\allowbreak a_{1}=3\) & \(m\le L\); & \(W_m T^L(F),\allowbreak  F=\prod_i h_{a_i} f_{b_i}\); & 2. \\
\addlinespace[2pt]
\(j\ge 2,\allowbreak a_{1}=3\) & \(m>L\); & \(\bar f_m AT^{L-1}(F),\allowbreak \) same \(F\); & 2. \\
\addlinespace[2pt]
\(a_{1}\ge 4,\allowbreak \text{any } j\ge 1\) & \(m\le L\); & \(W_m T^L(F),\allowbreak  F=\prod_i h_{a_i} f_{b_i}\); & 2. \\
\addlinespace[2pt]
\(a_{1}\ge 4,\allowbreak \text{any } j\ge 1\) & \(m>L\); & \(\bar f_m AT^{L-1}(F),\allowbreak \) same \(F\); & 2. \\
\end{longtable}
\endgroup


The bar in the minimum-one \(L=1\) row reverses every terminal residual
edge. The proof reverses the entire configuration,
including posts and pairs, when applying its analytic comparison;
the actual output in this table retains the prescribed negative \(P\).
Its condition \(m<u_{1}\) is preserved. Genuine equal pairs can either
follow that convention or be restored by \eqref{q3:eq:equalpair} and its complete reversal.

The first two rows use the center input of Section~\ref{proof:Q5-C113}.
For nonempty later-pair lists one may use rates $2005/192$ and $97/12$,
with $m<u_1$. Rows 3--8 follow Proposition~\ref{prop:R43:catalogue},
rows 9--12 follow Section~\ref{proof:R2-C038}, and rows 13--14
follow Section~\ref{proof:R2-C039}, with the directions in their
respective tables.

For \(j=1\), \(a_1\) is exactly \(1\), \(2\), \(3\), or at least \(4\). For \(j\ge2\), the same
four-way split applies. Thus every even list is covered.

\paragraph{R46.3.2 Odd terminal lists of minimum at least two}

Write \((\ell,\allowbreak r_1,\allowbreak t_1,\allowbreak \ldots,\allowbreak r_j,\allowbreak t_j)\), \(\ell\ge 2,\allowbreak  j\ge 1.\) Set

\fitdisplay{
F=\begin{cases}
 \bar f_\ell\prod_i h_{r_i}f_{t_i},&r_1\ge4,\\
 \bar f_2 f_3 h_{t_1}\prod_{i\ge2}h_{r_i}f_{t_i},&r_1=3.
 \end{cases}
}{R46.3.2}

The latter case forces \(\ell=2.\) Use

\fitdisplay{
H_0=\begin{cases}
 W_m T^L F,&m\le\ell+L,\\
 \bar f_m AT^{L-1}F,&m>\ell+L.
 \end{cases}
}{R46.3.3}

R2-C030/A014 gives eighth loss \(3\), including every permitted \(P\) and
arbitrary cross-block overlap. The comparison is with \(\ell+L\),
not merely \(L\). Every strict odd list with minimum at least \(2\) has second
entry \(3\) or at least \(4\), so (R46.3.2) exhausts it.

\paragraph{R46.3.3 Odd terminal lists of minimum one}

Write \((1,\allowbreak r_1,\allowbreak t_1,\allowbreak \ldots,\allowbreak r_j,\allowbreak t_j)\), \(j\ge 1.\) Choose


\begingroup\small
\setlength{\tabcolsep}{3pt}
\setlength{\LTleft}{0pt}\setlength{\LTright}{\fill}
\begin{longtable}{@{}p{0.2200\linewidth}p{0.7000\linewidth}@{}}
\caption{Ordinary eighth moments and the nonunit-pair outlet: parameter values.}\label{tab:parameters:r2-c046:2}\\
\toprule
Entry & Actual \(F\) \\
\midrule\endfirsthead
\toprule
Entry & Actual \(F\) \\
\midrule\endhead
\bottomrule\endfoot
\(r_{1}\ge 4\) & \(q \prod_i h_{r_i} f_{t_i}\). \\
\addlinespace[2pt]
\(r_{1}=3\) & \(q f_{3} h_{t_1} \prod_{i\ge 2}h_{r_i} f_{t_i}\). \\
\addlinespace[2pt]
\(r_{1}=2,\allowbreak j=1\) & \(q f_{2} h_{t_1}\). \\
\addlinespace[2pt]
exact \(j=2\) list \((1,2,3,4,e)\) & \(q f_{2} h_e h_{3} f_{4}\). \\
\addlinespace[2pt]
every other \(r_{1}=2\) & \(q f_{2} h_{t_1} \prod_{i\ge 2}h_{r_i} f_{t_i}\). \\
\end{longtable}
\endgroup


The actual anchor/connector is


\begingroup\small
\setlength{\tabcolsep}{3pt}
\setlength{\LTleft}{0pt}\setlength{\LTright}{\fill}
\begin{longtable}{@{}p{0.2200\linewidth}p{0.7000\linewidth}@{}}
\caption{Ordinary eighth moments and the nonunit-pair outlet: parameter values.}\label{tab:parameters:r2-c046:3}\\
\toprule
Entry & Unpaired anchor/connector factor \(H_0\) \\
\midrule\endfirsthead
\toprule
Entry & Unpaired anchor/connector factor \(H_0\) \\
\midrule\endhead
\bottomrule\endfoot
\(L=1,\allowbreak r_{1}=2\) & \(\bar f_m A F\). \\
\addlinespace[2pt]
\(L=1,\allowbreak r_{1}\ge 3\) & \(W_m T F\). \\
\addlinespace[2pt]
\(L\ge 2,\allowbreak m\le L+1\) & \(W_m T^L F\). \\
\addlinespace[2pt]
\(L\ge 2,\allowbreak m>L+1\) & \(\bar f_m AT^{L-1}F\). \\
\end{longtable}
\endgroup


Proposition~\ref{prop:r40:ordinary-catalogue} supplies ordinary moment bounds of orders \(8\) and \(4\)
for these same directions, with eighth loss
at least \(57/400\). The direction for \((1,\allowbreak 2,\allowbreak 3)\) is \(qf_2h_3\) even at \(L=2\);
it must not be replaced by C106's distinct ordinary-unit exception.
The crossed five-arm row uses exactly its existing arms.
Every unit-minimum odd list has second entry \(2\), \(3\), or at least \(4\);
the displayed cardinality/prefix splits are exhaustive.

\paragraph{R46.3.4 Proof of the common theorem and exact post scope}

Each cited local theorem holds on all original finite hosts,
all integers in its row and the stated complete finite post block.
All its displayed eighth losses are at least \(1/1024.\) Genuine
equal pairs only decrease the same positive message by \eqref{q3:eq:equalpair} and its complete reversal.
This proves (R46.1.2) for the complete union.

In every later
central application, the anchor is selected before all post
lengths in the strict central list, so \(m<u_{1}\) and \(u_{1}\ge 4\) hold
automatically. At a nonunit-pair root there is no unequal post.


This theorem has \(m\ge3\). It gives no general eighth-moment bound
for \(m=2\) and no standalone \(M=2\) outlet for a singleton residual
by itself.

\paragraph{R46.4 Every nonunit pair, all inner and outer lengths}

At an actual internal root with strict residual \((s,m)\), \(2\le s<m\),
orient its \(s\)-arm inward as \(\bar f_s\). Put its distinct genuine \(m\)-arm,
terminal, \(L\) connector and all its root/terminal equal pairs in
the actual \(H\) of \mbox{\hyperref[part:R46.3]{part R46.3}}, with \(P=1\). Then the exact outgoing
message along an outer connector of length \(M\ge 1\) directed toward
the receiving center is

\fitdisplay{
U=T^M(\bar f_s H).
}{R46.4.1}

The \(m\)-arm has length at least \(3\). The \(s\)-arm is a different original arm,
even when its length coincides with a length in the other
terminal's table. Each actual edge is used once.

The full-multiplier inequality of Section~\ref{proof:Q5-C066} holds for every real \(Y\):

\fitdisplay{
I_q(T^M(\bar f_aY))^3\le(1-d/72)E|Y|^8,\qquad a\ge2,M\ge1.
}{R46.4.2}

The multiplier \(T^M M_{\bar f_a}\) retains the inward arm in this
operator order; its proof uses the row \(\bar f_{M+a}\). Thus \(Y\)
may be the \(AT^{L-1}\) component of the preceding directory.

Apply (R46.4.2) with \(a=s,\allowbreak  Y=H\) and use (R46.1.2):

\fitdisplay{
I_q(U)^3\le(1-d/72)e^{-d/1024}
          \le e^{-137d/9216}.
}{R46.4.3}

The full multiplier includes every \(M\ge1\), independently of \(L\).
The retained inward arm is the existing \(s\)-arm, distinct from the
\(m\)-arm used in \(H\).

The path construction combines this local output with the outlet
on the other side using \eqref{eq:five-chain-close}. The tournament
and entropy consequences are proved in Section~\ref{sec:completion}.
